\section*{Введение}
\addcontentsline{toc}{section}{Введение}

\textbf{Актуальность темы.} В условиях активного развития фитнес-индустрии
и роста популярности здорового образа жизни всё большее число людей обращается
к услугам фитнес-центров. Управление записью клиентов, расписанием тренировок
и взаимодействием с тренерами вручную становится неэффективным и приводит к
потере клиентов и снижению качества сервиса. Разработка специализированных
веб-приложений, автоматизирующих эти процессы, является актуальной задачей
для современного фитнес-бизнеса.

\textbf{Степень разработанности проблемы.} На рынке существуют коммерческие
решения для автоматизации фитнес-центров (напр. Mindbody, 1С:Фитнес клуб),
однако они ориентированы на крупные сети и обладают избыточным функционалом
для небольших студий. Разработка доступного и гибкого веб-приложения с открытым
исходным кодом остаётся практически востребованной задачей.

\textbf{Объект исследования} — процессы управления записью клиентов,
расписанием тренировок и взаимодействием пользователей в фитнес-центре.

\textbf{Предмет исследования} — программные средства для автоматизации
онлайн-записи на тренировки и управления личным кабинетом пользователя
в рамках веб-ориентированной платформы.

\textbf{Целью данной работы} является проектирование и реализация
веб-приложения для фитнес-центра, обеспечивающего онлайн-запись на тренировки,
управление расписанием и личный кабинет пользователя с разграничением прав
доступа по ролям.